\documentclass[a4paper,10pt]{article}

\usepackage[
    a4paper,
    top=1.0cm,
    bottom=0.9cm,
    left=1.45cm,
    right=1.35cm
]{geometry}

\usepackage{fontspec}
\usepackage{polyglossia}
\setdefaultlanguage[babelshorthands=true]{russian}
\setotherlanguage{english}

\setmainfont{FreeSerif}
\setmonofont{DejaVuSansMono}

\usepackage{microtype}
\usepackage{parskip}
\setlength{\parindent}{0pt}
\setlength{\parskip}{0.08em}

\usepackage{titlesec}
\titlespacing*{\section}{0pt}{0.35em}{0.10em}

\usepackage{xurl}
\usepackage{hyperref}
\usepackage{bookmark}
\usepackage{array}
\usepackage{booktabs}
\usepackage{float}

\hypersetup{
    colorlinks=true,
    linkcolor=black,
    urlcolor=blue,
    citecolor=black,
    breaklinks=true
}

\urlstyle{same}

\usepackage{titling}
\setlength{\droptitle}{-0.65cm}

\begin{document}

\raggedbottom

\title{Анализ графов}
\author{Пентегова Мария Сергеевна}
\date{}
\maketitle

\begin{flushright}
Группа: 25.Б41-мм \\
Кафедра: системного программирования \\
Научный руководитель: Григорьев Семён Вячеславович \\
Номер семестра практики: 2
\end{flushright}

\section{Введение и постановка задачи}

Работа выполнена в рамках учебной практики по анализу графов с
использованием GraphBLAS --- программного интерфейса для представления
алгоритмов над графами в терминах операций линейной алгебры над
разреженными матрицами и векторами~\cite{graphblas-api}. В работе
используется SuiteSparse:GraphBLAS --- реализация GraphBLAS на языке
C~\cite{graphblas}.

Рассматриваемый алгоритм --- обход в ширину (Breadth-First Search,
BFS), при котором вершины посещаются по уровням расстояния от исходных
вершин. В работе реализованы два варианта такого обхода: BFS с
построением массива родителей вершин (parent-BFS) и BFS из нескольких
исходных вершин (multisource BFS).

\textbf{Цель работы} --- реализовать варианты BFS на языке C в
классическом представлении графа и с использованием SuiteSparse:GraphBLAS,
а также сравнить время их выполнения на графах из SuiteSparse Matrix
Collection~\cite{ssmc}.

Для достижения цели поставлены следующие задачи:
\begin{enumerate}
    \item реализовать parent-BFS и multisource BFS с использованием
    представления графа в формате CSR;
    \item реализовать соответствующие варианты BFS средствами
    SuiteSparse:GraphBLAS;
    \item проверить корректность реализованных алгоритмов;
    \item измерить время выполнения на семи графах из SuiteSparse Matrix
    Collection;
    \item сопоставить результаты классического и GraphBLAS-подходов и
    определить, в каких исследованных случаях каждый из них имеет
    меньшее измеренное время выполнения.
\end{enumerate}

\section{Реализация}

В классических вариантах используется CSR (Compressed Sparse Row) ---
формат хранения разреженной матрицы, в котором номера соседей каждой
вершины располагаются последовательно. Для вершины $v$ её соседи
находятся в диапазоне
\[
[\texttt{row\_ptr}[v],\texttt{row\_ptr}[v+1])
\]
массива \texttt{col\_idx}. При первом посещении вершины $u$ из вершины
$v$ parent-BFS сохраняет
\[
\texttt{parent}[u]=v.
\]
В multisource BFS очередь изначально содержит несколько исходных
вершин. Для CSR-представления один запуск BFS имеет асимптотическую
сложность $O(|V|+|E|)$.

Для сравнения могли быть использованы LAGraph и
Boost Graph Library (BGL). LAGraph предоставляет алгоритмы анализа
графов поверх GraphBLAS, а BGL содержит обобщённые структуры данных и
алгоритмы работы с графами~\cite{lagraph,boost-graph}. Однако готовые
реализации BFS из этих библиотек в измеряемых вариантах не
использовались.

Самостоятельная реализация необходима для контролируемого сравнения
двух способов выполнения одного алгоритма. В эксперименте классический
BFS и BFS на основе GraphBLAS используют одинаковые входные графы,
исходные вершины и условия запуска. Использование готового алгоритма
из LAGraph или BGL не позволило бы рассматривать его внутреннюю
реализацию как часть контролируемого варианта сравнения. Поэтому
классические варианты были реализованы самостоятельно, а GraphBLAS
варианты непосредственно используют API SuiteSparse:GraphBLAS.

В GraphBLAS граф представляется булевой матрицей смежности $A$, а
текущий фронт обхода --- разреженным булевым вектором $f_k$. Следующий
фронт формируется матрично-векторной операцией
\[
f_{k+1}=f_k A.
\]
Уже посещённые вершины исключаются с помощью маски. Номер итерации
соответствует уровню вершины в BFS. В multisource-варианте начальный
вектор содержит несколько исходных вершин.

Исходные графы загружаются из файлов Matrix Market --- формата
представления разреженных матриц~\cite{ssmc}. Матрицы с ключевым словом
\texttt{general} в заголовке рассматриваются как ориентированные
графы, а матрицы с маркировкой \texttt{symmetric} --- как
неориентированные. Перед выполнением BFS направление рёбер не
изменяется: происходит только преобразование структуры хранения
исходных списков смежности в формат CSR. Для ориентированных графов
обход выполняется по исходящим связям.

При выборе начальной вершины для однократного BFS используется вершина
с максимальной исходящей степенью (out-degree), то есть числом дуг,
выходящих из неё. Для multisource BFS используются четыре источника
с индексами $0$, $n/4$, $n/2$ и $3n/4$.

\section{Эксперимент}

Цель эксперимента --- сравнить время выполнения четырёх реализаций BFS
на графах различного размера и структуры. В эксперимент включены семь
графов из SuiteSparse Matrix Collection~\cite{ssmc}.

Значения $|V|$ и $|E|$ считываются непосредственно из заголовка и
координатных записей соответствующего файла Matrix Market. В
эксперименте $|V|$ соответствует размеру квадратной матрицы смежности,
а $|E|$ --- числу координатных записей. Для ориентированных матриц с
маркировкой \texttt{general} каждая такая запись соответствует одной
дуге. Для матриц с маркировкой \texttt{symmetric} используется
неориентированная интерпретация графа.

\begin{table}[H]
\centering
\scriptsize
\setlength{\tabcolsep}{3.5pt}
\begin{tabular}{lrr}
\toprule
Граф & $|V|$, млн & $|E|$, млн \\
\midrule
cit-Patents & 23.947 & 28.854 \\
com-Amazon & 0.335 & 0.926 \\
com-Orkut & 3.072 & 117.185 \\
roadNet-CA & 0.916 & 5.105 \\
road\_usa & 3.775 & 16.519 \\
soc-LiveJournal1 & 4.848 & 68.994 \\
web-Google & 1.971 & 2.767 \\
\bottomrule
\end{tabular}
\caption{Характеристики графов, использованных в эксперименте.}
\label{tab:graphs}
\end{table}

Эксперимент выполнен на процессоре Intel Core Ultra 7 255H с 30~ГБ
оперативной памяти под управлением Ubuntu 26.04.1 LTS. Использованы
GCC 15.2.0 и SuiteSparse:GraphBLAS 7.4.0. Программа собирается со
стандартом C11 с параметрами \texttt{-O3 -DNDEBUG}. Параметр
\texttt{-O3} включает высокий уровень автоматической оптимизации
компилятора, а \texttt{-DNDEBUG} отключает проверки \texttt{assert}
во время измерений.

Для каждой реализации на каждом графе выполняется 10 независимых
измерений времени выполнения BFS. В качестве основной метрики
используется среднее арифметическое полученных значений. Минимальное
и максимальное время также вычисляются программой, но не включены в
итоговую таблицу, поскольку основной целью сравнения является
сопоставление средних значений.

Построение матрицы смежности GraphBLAS выполняется до серии измерений
и не включается в измеряемое время. Таким образом, сравнивается именно
время выполнения алгоритмов BFS при уже подготовленном представлении
графа.


\begin{table}[H]
\centering
\scriptsize
\setlength{\tabcolsep}{2.5pt}
\begin{tabular}{lrrrr}
\toprule
Граф & CP, мс & CM, мс & GB-L, мс & GB-M, мс \\
\midrule
cit-Patents & 1034.142 & 1160.571 & 26347.954 & 24488.670 \\
com-Amazon & 29.881 & 22.404 & 62.855 & 61.161 \\
com-Orkut & 2121.789 & 1993.293 & 542.795 & 582.770 \\
roadNet-CA & 52.984 & 48.788 & 97.838 & 109.488 \\
road\_usa & 7.797 & 0.491 & 189.076 & 36.194 \\
soc-LiveJournal1 & 1613.440 & 1662.308 & 439.515 & 390.446 \\
web-Google & 88.148 & 79.361 & 713.218 & 688.762 \\
\bottomrule
\end{tabular}
\caption{Среднее время выполнения BFS по 10 запускам, мс.}
\label{tab:results}
\end{table}

В таблице CP обозначает классический parent-BFS, CM --- классический
multisource BFS, GB-L --- GraphBLAS BFS с вычислением уровней, GB-M ---
GraphBLAS multisource BFS.

\section{Анализ результатов}

Проведённые измерения показывают, что производительность
классического CSR-подхода и GraphBLAS различается в зависимости от
исследуемого графа.

Для parent-BFS GraphBLAS-вариант показал меньшее среднее время на
графах \texttt{com-Orkut} и \texttt{soc-LiveJournal1}. На
\texttt{com-Orkut} среднее время GraphBLAS Level BFS составило
542.795~мс против 2121.789~мс у классического parent-BFS, а на
\texttt{soc-LiveJournal1} --- 439.515~мс против 1613.440~мс.
На остальных исследованных графах меньшее среднее время показал
классический parent-BFS.

Для multisource BFS получена аналогичная картина. GraphBLAS
Multisource BFS показал меньшее среднее время на
\texttt{com-Orkut} и \texttt{soc-LiveJournal1}: 582.770~мс против
1993.293~мс и 390.446~мс против 1662.308~мс соответственно. На
остальных пяти графах меньшее среднее время показал классический
multisource BFS.

Таким образом, в проведённом эксперименте GraphBLAS продемонстрировал
преимущество на крупномасштабных социальных графах, представленных
\texttt{com-Orkut} и \texttt{soc-LiveJournal1}. На остальных
исследованных структурах меньшее время выполнения показали
классические реализации на основе CSR. Это позволяет сделать
практический вывод о том, что выбор представления и способа
реализации BFS может существенно влиять на время обработки больших
разреженных графов.

При этом эксперимент не включал отдельное измерение размеров фронтов
BFS, времени обработки отдельных уровней и других топологических
характеристик. Поэтому конкретные причины наблюдаемого различия в
производительности в данной работе не устанавливаются. Зафиксированное
различие во времени является непосредственно измеренным результатом
эксперимента.


\section{Проверка корректности}

Корректность реализованных вариантов BFS проверяется
автоматизированным тестовым набором. Проверка выполняется как для
отдельных реализаций, так и посредством перекрёстного сравнения
классического и GraphBLAS-подходов.

Для отдельных реализаций проверяются следующие свойства:
\begin{itemize}
    \item достижимость всех доступных вершин на соответствующих
    тестовых графах;
    \item корректность вычисленных уровней BFS;
    \item корректность массива родителей для parent-BFS;
    \item отсутствие циклов в цепочках родителей и связь посещённых
    вершин с исходной вершиной.
\end{itemize}

Изолированные тесты выполняются для классического parent-BFS,
классического multisource BFS, GraphBLAS BFS и GraphBLAS
multisource BFS. Тестовые функции используют заранее заданные
топологии, для которых ожидаемые свойства результата могут быть
проверены автоматически.

Дополнительно выполняется перекрёстная проверка результатов
классических и GraphBLAS-реализаций. Для single-source BFS
классическая реализация на основе CSR и GraphBLAS-реализация
запускаются на одинаковом тестовом графе и из одной исходной
вершины, после чего выполняется поэлементное сравнение полученных
векторов уровней. Для multisource BFS аналогичная проверка
выполняется при одинаковом наборе исходных вершин.

Такое сравнение позволяет проверить не только внутренние инварианты
отдельных реализаций, но и соответствие результатов двух
принципиально разных способов реализации BFS: явного обхода CSR и
матрично-векторного обхода средствами GraphBLAS.

Автоматический тестовый набор завершился результатом
\texttt{100\% passed}. Это означает, что все предусмотренные
тестовые случаи успешно прошли заданные проверки. Полученный
результат относится к использованным тестовым топологиям и
проверяемым сценариям и не является формальным доказательством
корректности для любого возможного графа.

\section{Заключение}

В ходе учебной практики реализованы и экспериментально сопоставлены
четыре варианта BFS на языке C: классический parent-BFS, классический
multisource BFS, GraphBLAS Level BFS и GraphBLAS Multisource BFS.
Классические реализации используют CSR-представление графа, тогда как
GraphBLAS-варианты используют булеву матрицу смежности и
матрично-векторные операции над разреженными структурами.

Проведено экспериментальное сравнение четырёх реализаций на семи
графах из SuiteSparse Matrix Collection. Для каждого варианта на каждом
графе выполнено по 10 измерений времени, а для сравнения использовалось
среднее значение. Корректность реализаций проверена автоматизированным
тестовым набором, включающим проверку достижимости, уровней BFS,
массива родителей и соответствия результатов классического и
GraphBLAS-подходов. Все предусмотренные тесты успешно пройдены
(\texttt{100\% passed}).

Эксперимент показал, что производительность BFS существенно зависит
от структуры входного графа. GraphBLAS-варианты показали меньшее
среднее время на двух исследованных крупномасштабных социальных
графах, тогда как на остальных пяти графах меньшее время выполнения
показали классические реализации на основе CSR. Таким образом,
поставленная задача сравнения двух способов реализации BFS выполнена
не только на уровне алгоритмической реализации, но и
экспериментально: получены измерения, позволяющие определить различие
в производительности подходов на графах различной структуры.

Результаты работы показывают, что выбор способа реализации BFS
оказывает заметное влияние на время выполнения даже при использовании
одного и того же алгоритма. Проведённое сравнение даёт практическую
основу для выбора реализации при работе с графами, близкими по
структуре к исследованным в эксперименте.

Дальнейшее развитие работы может включать измерение размеров фронтов
BFS на каждом уровне, времени обработки отдельных уровней и других
характеристик обхода. Это позволит дополнительно исследовать причины
наблюдаемого различия производительности между CSR- и
GraphBLAS-реализациями.


\section*{Репозиторий проекта}

Исходный код, тесты, конфигурация сборки и материалы эксперимента
размещены в репозитории проекта~\cite{repository}. Файлы матриц
формата Matrix Market хранятся с использованием Git Large File Storage
(Git LFS) --- расширения Git для хранения крупных файлов~\cite{git-lfs}.

\begin{thebibliography}{8}

\setlength{\itemsep}{0.05em}
\setlength{\parskip}{0pt}

\bibitem{graphblas-api}
GraphBLAS Forum.
\textit{The GraphBLAS C API Specification. Version 2.1}, 2023.
URL: \url{https://graphblas.org/docs/GraphBLAS_API_C_v2.1.0.pdf}
(дата обращения: 24.09.2026).

\bibitem{graphblas}
Davis, T. A.
\textit{SuiteSparse:GraphBLAS}.
Официальный репозиторий проекта.
URL: \url{https://github.com/DrTimothyAldenDavis/GraphBLAS}
(дата обращения: 24.09.2026).

\bibitem{lagraph}
GraphBLAS.org.
\textit{LAGraph}.
Официальный репозиторий проекта.
URL: \url{https://github.com/GraphBLAS/LAGraph}
(дата обращения: 24.09.2026).

\bibitem{boost-graph}
Boost.
\textit{Boost Graph Library}.
Официальная документация.
URL: \url{https://www.boost.org/library/latest/graph/}
(дата обращения: 24.09.2026).

\bibitem{ssmc}
Davis, T. A., Hu, Y.
\textit{The University of Florida Sparse Matrix Collection}.
ACM Transactions on Mathematical Software, 38(1), 2011.
URL: \url{https://sparse.tamu.edu/}
(дата обращения: 24.09.2026).

\bibitem{cmake-ctest}
Kitware.
\textit{CTest Documentation}.
URL: \url{https://cmake.org/cmake/help/latest/manual/ctest.1.html}
(дата обращения: 24.09.2026).

\bibitem{git-lfs}
Git LFS.
\textit{Git Large File Storage}.
URL: \url{https://git-lfs.com/}
(дата обращения: 24.09.2026).

\bibitem{repository}
Pentegova, M.
\textit{graph\_blas\_bfs-2}.
Репозиторий проекта.
URL: \url{https://github.com/MariaPentegova/graph_blas_bfs-2}
(дата обращения: 24.09.2026).

\end{thebibliography}

\end{document}

